\documentclass[preprint,12pt]{elsarticle}
% Alternative Optionen:
% \documentclass[final,3p,times]{elsarticle}         % 3-spaltig final
% \documentclass[final,5p,times,twocolumn]{elsarticle} % 2-spaltig final

% --- Pakete ---
\usepackage{amsmath,amssymb}
\usepackage{graphicx}
\usepackage{lineno}          % Zeilennummerierung (oft gefordert)
\usepackage{booktabs}        % schöne Tabellen
\usepackage{siunitx}         % SI-Einheiten
\usepackage[hidelinks]{hyperref}

% Zitierstil (Elsevier: numerisch)
\biboptions{sort&compress}

\journal{Energy Conversion and Management}

\begin{document}

\begin{frontmatter}

\title{Titel der Publikation}

%% Autoren
\author[inst1]{Vorname Nachname\corref{cor1}}
\ead{email@institut.de}
\author[inst2]{Zweiter Autor}

\cortext[cor1]{Corresponding author}

\affiliation[inst1]{organization={Institut / Abteilung},
            addressline={Straße},
            city={Stadt},
            postcode={PLZ},
            country={Land}}
\affiliation[inst2]{organization={Zweite Einrichtung},
            city={Stadt},
            country={Land}}

%% Abstract
\begin{abstract}
Hier steht die Kurzfassung (ca. 150--250 Wörter).
\end{abstract}

%% Highlights (separat als Bullet-Points, 3--5 Stück, je <85 Zeichen)
\begin{highlights}
\item Erster Hauptbeitrag.
\item Zweiter Hauptbeitrag.
\item Dritter Hauptbeitrag.
\end{highlights}

%% Keywords
\begin{keyword}
Keyword1 \sep Keyword2 \sep Keyword3
\end{keyword}

\end{frontmatter}

\linenumbers   % Zeilennummerierung aktivieren

%% ============ HAUPTTEXT ============
\section{Introduction}
\label{sec:intro}
Einleitung...

\section{Methodology}
\label{sec:method}
Methoden...

\section{Results and Discussion}
\label{sec:results}
Ergebnisse...

\section{Conclusions}
Schlussfolgerungen...

%% Nomenklatur (in ECM häufig verwendet)
% \section*{Nomenclature}

%% Danksagung
\section*{Acknowledgements}
Dank / Fördermittel.

%% ============ REFERENZEN ============
% Variante A: BibTeX mit Elsevier-Stil
\bibliographystyle{elsarticle-num}   % oder elsarticle-harv
\bibliography{references}

% Variante B: manuell
% \begin{thebibliography}{00}
% \bibitem{ref1} Autor. Titel. Journal Jahr;Vol:Seiten.
% \end{thebibliography}

\end{document}